\section{Introduction}



\begin{figure*}[t]
  \includegraphics[width=\textwidth]{figs/task_compare_v3.png}
  \caption{Input and output example of the TVS task as a side-by side comparison to a superficially similar task: temporal localization. TVS focuses on goal-driven cognitive alignment and performs reasoning-guided problem reconstruction through information screening, whereas grounding only emphasizes perceptual alignment and locates depictive level visual evidence by direct moment retrieval.\vspace{-10pt}}
  \label{fig:tvs_task_compare}
\end{figure*}

Human cognition optimizes information processing via a two-stage control mechanism: a fast, pre-attentive screening to prune redundancy, followed by focused reasoning on a compacted, goal-aligned representation \citep{FITNO1, GS2.0NO2, zacks2007event, topdownfeedback}. Behaviorally, this manifests as scrubbing the seek bar before detailed viewing \citep{wu2024v}, minimizing Extraneous Load to purify Germane Load \citep{cognitiveLoadDuringProblemSolving:EffectsOnLearningNO3, NineWaysToReduceCognitiveLoadInMultimediaLearningNO4, ReducingExtraneousLoadImportantNO15}---reshaping the problem before engaging the expensive reasoning stage.

This filter-before-reason economy is forced by VideoQA: a minute of video contains thousands of frames of which only a narrow band bears on any given query, yet current Video-LLMs \citep{chen2023video-LLMmodelingvideosequence,2023videochat,Xu2024PLLaVAP,maaz-etal-2024-video,li2024mvbenchcomprehensivemultimodalvideo,Ma_2024_CVPR} typically fit the entire stream through a bounded visual-token budget via query-agnostic uniform sampling. This dilutes semantic signal at critical segments and injects noise, weakening supervision during training \citep{lin2024norton, knowledgeovershadowing} and reducing task-oriented focus at inference.

Existing work such as traditional grounding and grounded QA \citep{nextgqa, cgbench} attempts the obvious remedy of retrieving query-relevant frames, but operates at the \textbf{perceptual level}---depictive similarity between query tokens and visual evidence. Complex queries entangle multi-hop temporal, relational, and causal dependencies with sparse, weakly localized cues, so retrieving depicted frames without adapting the query leaves a multimodal semantic mismatch that misleads the downstream model.

This calls for a novel mechanism that jointly reshapes reasoning structure and cognitive load while treating input modalities \textbf{as a whole}. The two-stage structure of human cognition \citep{FITNO1, GS2.0NO2} suggests a clean architectural response: decouple screening from reasoning so one mechanism serves arbitrary downstream reasoners. We design a \textbf{front-end adapter} that reshapes the raw instance symmetrically over all modalities before it reaches any Video-LLM.

For such an adapter to compose with any downstream agent, its output must share its input's signature, forcing the transformation to be \textbf{endomorphic}:
\[
\mathcal{F}_{\text{TVS}} : (V, Q) \rightarrow (v, q),
\]
mapping a multimodal task instance back into the same task space. This endomorphism is the critical enabling condition for model-agnostic answer invariance---that any reasonable agent $\mathcal{M}$ yield the same answer on $(v, q)$ as on $(V, Q)$---which is only well-defined when compressed and original inputs share the same model signature. The formulation then extends naturally to \textbf{any multimodal task} admitting a uniform agent signature; VideoQA is one instance.

Taking a generative view---where answer $A$ is a latent fact and $(V, Q)$ is an observation drawn from it---this compression forms the Markov chain $A \to (V, Q) \to (v, q)$, and \underline{\textbf{T}}emporal \underline{\textbf{V}}isual \underline{\textbf{S}}creening (TVS) is precisely a \textbf{structurally-constrained sufficient-statistic problem}: find the most compact admissible $(v, q)$ preserving all $A$-relevant information. This identification places TVS in the Information Bottleneck lineage \citep{tishby1999information, tishby2015deep}, instantiated over the original multimodal input space rather than a learned latent code. $A$ is conserved throughout, aligning with the fixity of the Goal State in Problem Space Theory \citep{HumanProblemSolvingNO12}: TVS reshapes ``what to ask'', not ``what to answer''.

TVS yields four concrete benefits: (i) Training alignment: removing off-target supervision boosts gradient signal-to-noise and reduces spurious correlations under fixed budgets; (ii) Inference robustness: offloading non-informative perceptual burdens shortens the implicit reasoning program, unleashing reasoning capabilities; (iii) Interpretability: the compact $(v, q)$ is a controllable, faithful artifact exposing model bottlenecks; (iv) Generalizability: as a model-agnostic front-end, TVS benefits diverse VideoQA agents.

We present \textit{ReSimplifyIt}, a plug-and-play multi-agent framework that \underline{\textbf{ReS}}iliently \underline{\textbf{imp}}rovises and qua\underline{\textbf{lif}}ies proposals \underline{\textbf{It}}eratively. Each screening round runs a language-only Launcher that hypothesizes a trimming instruction and rewritten query (using the intrinsic cross-modal relation as a prior), a Validator that executes and critiques these plans through a Viewer module, and memory trackers for self-correction---operationalizing a competence-from-consequence~\citep{brooks1991intelligence} principle: trial execution provides constructive feedback tightening the coupling between hypothesized reasoning structure and available evidence.

We further construct YouCookII-TVS, the first benchmark dedicated to evaluating TVS as a joint transformation over video and query. Benchmarking shows TVS is a non-trivial task far from solved, revealing a key bottleneck for downstream VideoQA reasoning and deserving independent study.

To our knowledge, ours is the first work to formally recognize that grounding and trimming a video creates a semantic mismatch with complex, context-dependent queries, to formalize the task on this premise, and to extensively explore temporal filtering within the Video-LLM framework. Our results show that TVS not only improves inference-time performance but also enables more structured and effective training through finer-grained multimodal alignment, paving the way for scalable, interpretable, cognitively-inspired video-language understanding systems.

The main contributions of this work are summarized below:

\begin{itemize}[leftmargin=*,topsep=-1pt,itemsep=0ex]
\item We propose temporal visual screening (TVS), a novel cognitively-inspired, information-theoretically framed task that guides task-aware reasoning structure filtering and addresses a key bottleneck in multi-modal alignment. We also release YouCookII-TVS, a benchmark with 238.2 hours of video and 2{,}754 questions.
\item We introduce ReSimplifyIt, the first baseline for the TVS task: a model-agnostic plug-in compatible with \textbf{any} existing videoQA or instruction-tuning pipeline.
\item We quantitatively evaluate TVS in VideoQA and Video Instruction Tuning, where our method delivers over $10\%$ and $5\%$ absolute gains at inference and training, respectively, across multiple strong baselines and benchmarks.
\end{itemize}

\begin{figure*}[t]
  \includegraphics[width=\textwidth]{figs/workflow_figure_new_eg-3_manlingstyle.png}  
  \caption{Snapshot examples of the workflow of our proposed ReSimplifyIt framework.}
  \label{fig:workflow_figure}
\end{figure*}
